\documentclass[12pt, a4paper]{article}
\usepackage{graphicx}
\usepackage{amsmath}
\usepackage{amssymb}
\usepackage{listings}
\usepackage{xcolor}

% Configuración de estilo para el código Scala
\lstdefinestyle{scalaStyle}{
    language=Scala,                 % Lenguaje del código
    basicstyle=\ttfamily\small,     % Estilo básico de fuente y tamaño
    keywordstyle=\color{blue},      % Estilo para palabras clave
    commentstyle=\color{green},     % Estilo para comentarios
    numbers=left,                   % Mostrar números de línea a la izquierda
    numberstyle=\tiny,              % Estilo de los números de línea
    frame=single,                   % Un marco alrededor del código
    breaklines=true,                % Permitir saltos de línea automáticos
    captionpos=b,                   % Posición de la leyenda
}

\graphicspath{ {./images/} }
\title{
  \begin{figure}[th]
    \centering
    \includegraphics[width=0.2\textwidth]{Univalle}
  \end{figure}
  \textbf{Universidad del Valle
    \\{\Large Facultad de ingeniería}
  \\{\large Ingeniería en sistemas}}}
\author{Cristian David Pacheco Torres
  \\ 2227437
  \\ Juan Sebastian Molina Cuellar
  \\ 2224491}
\date{Septiembre 2023}
\begin{document}
\maketitle
Taller 2
\newpage{}
\begin{abstract}
Your abstract goes here functional programming
\end{abstract}
\newpage{}
\tableofcontents
\newpage{}
\section{Introduction}
A introduction a ver asdf 
\section{Taller 1 : Recursión}
Para el desarrollo de este taller, se utilizaron las siguientes funciones en scala: \\
\begin{itemize}
  \item $isEmpty$: Boolean (Devuelve si una lista esta vacia).
  \item $head$: Int (devuelve si una lista $l$ esta vacia).
  \item $tail$: List[Int] (devuelve la lista sin el primer elemento $l$).
  \item $x::l$: devuelve la lista que representa la secuencia $<x,x_1,x_2,...,x_n>$ si $l$ es la lista que representa
  la secuencia $<x_1,x_2,...,x_n>$.
  \item $l1 ++ l2$ devuelve la lista que representa la concatenacion de las secuencias representadas por $l1$ y $l2$.
\end{itemize}
Se da por hecho de que las funciones anteriores estan argumentadas y demostradas, por lo tanto se procede a la resolucion de los ejercicios.
\subsection{Calcular el tamaño de una lista con un proceso iterativo}
En el listing 1 se puede apreciar la funcion $tamR(l)$ que a traves de un proceso recursivo calcula el tamaño de una lista.
\begin{lstlisting}[style=scalaStyle, caption=Calcula el tamaño de una lista con un proceso recursivo]
  def tamR(l : List[Int] ) : Int = {
    if (l.isEmpty) 0
    else 1 + tamR(l.tail )
  }
\end{lstlisting}
El problema a solucionar es hacer una funcion $tamI(l)$ tal que: \\
$tamR(l) == tamI(l)$
\subsubsection{Informe de procesos}
\textbf{Descripcion de la funcion.}\\
La solucion propuesta se basa en el siguiente algoritmo (ver listing 2.) \\
Para ello se creo la funcion $tamI(l:List[Int]):Int = \{\}$.
\begin{lstlisting}[style=scalaStyle, caption=Calcula el tamaño de una lista con un proceso iterativo]
  def tamI(l : List[Int]): Int = {
    def tam(lst : List[Int], acc : Int): Int = {
      if (lst.isEmpty) acc
      else tam(lst.tail, acc + 1)
    }    
  tam(l, 0)
  } 
\end{lstlisting}
Para esta implementacion se utilizo una funcion auxiliar: \\ $tam(lst:List[Int], acc:Int):Int$ (ver linea 2) la cual recibe como parametros una lista ($l$) y un acumulador ($acc$), el cual se encarga de contar el tamaño de la lista.
En el caso base (ver linea 3), si la lista esta vacia, se retorna el acumulador, en caso contrario se llama a la funcion $tam$ (ver linea 4) con la lista sin el primer elemento ($lst.tail$) y el acumulador incrementado en 1. \\
Para la inicializacion de la funcion $tamI(l:List[Int]):Int$ se llama a la funcion $tam$ (linea 6) con la lista y un acumulador inicializado en 0.\\ \\
\textbf{Tipo de proceso.}\\
Se pretende de que el tipo de proceso es \textbf{iterativo} para ello evaluamos la funcion con parametro $List(12, 3, 1, 8, 4)$ \\
$tamI(List(12, 3, 1, 8, 4))$\\
$\rightarrow  tamI(List(12, 3, 1, 8, 4))$ \\
$\rightarrow tam(List(12,3,1,8,4),0)$\\
$\rightarrow if(List(12,3,1,8,4).isEmpty) ~0 ~else ~tam(List(12,3,1,8,4).tail,0+1)$ \\
$\twoheadrightarrow tam(List(3,1,8,4),1)$\\
$\rightarrow if(List(3,1,8,4).isEmpty) ~1 ~else ~tam(List(3,1,8,4).tail,1+1)$ \\
$\twoheadrightarrow tam(List(1,8,4),2)$\\
$\rightarrow if(List(1,8,4).isEmpty) ~2 ~else ~tam(List(1,8,4).tail,2+1)$ \\
$\twoheadrightarrow tam(List(8,4),3)$\\
$\rightarrow if(List(8,4).isEmpty) ~3 ~else ~tam(List(8,4).tail,3+1)$ \\
$\twoheadrightarrow tam(List(4),4)$\\
$\rightarrow if(List(4).isEmpty) ~4 ~else ~tam(List(4).tail,4+1)$ \\
$\twoheadrightarrow tam(List(),5)$\\
$\rightarrow if(List().isEmpty) ~5 ~else ~tam(List().tail,5+1)$ \\
$\twoheadrightarrow 5$\\
Puesto que el proceso mantiene una forma constante, se puede decir que el proceso es \textbf{iterativo}.
\subsubsection{Informe de corrección}
\subsection{Dividiendo una lista en dos sublistas a partir de un pivote}
\subsubsection{Informe de procesos}
\begin{lstlisting}[style=scalaStyle, caption=Calcula una lista construida con los valores menores de un valor pivote proporcionado]
   def noMenoresQue(l: List[Int], value : Int): List[Int] = {

    if (l.isEmpty) Nil
    else if (l.head >= value) l.head :: noMenoresQue(l.tail, value)
    else noMenoresQue(l.tail, value)
  }
\end{lstlisting}
Este es un proceso de forma expansion contracion de tiempo lineal en funcion  y es una funcion recursiva
$menoresQue(List(9, ~3, ~1, ~8, ~7, ~10, ~0, ~7), ~8)$\\
$\rightarrow if ~(List(9, ~3, ~1, ~8, ~7, ~10, ~0, ~7).isEmpty())~[~] \\
\ else~ if(9 < 8) ~9::List(3, ~1, ~8, ~7, ~10, ~0, ~7) \\
  ~else ~menoresQue(List(3, ~1, ~8, ~7, ~10, ~0, ~7), ~8)$ \\
$\twoheadrightarrow 3::menoresQue(List(1, ~8, ~7, ~10, ~0, ~7), ~8)$\\
$\twoheadrightarrow 3::1::menoresQue(List( 8, ~7, ~10, ~0, ~7), ~8)$ \\
$\twoheadrightarrow 3::1::menoresQue(List( 7, ~10, ~0, ~7), ~8)$ \\
$\twoheadrightarrow 3::1::7::menoresQue(List( 10, ~0, ~7), ~8)$ \\
$\twoheadrightarrow 3::1::7::menoresQue(List( 0, ~7), ~8)$ \\
$\twoheadrightarrow 3::1::7::0::menoresQue(List( 7 ), ~8)$ \\
$\twoheadrightarrow 3::1::7::0::7::menoresQue([~], ~8)$ \\
$\rightarrow 3::1::7::0::[7]$ \\
$\rightarrow 3::1::7::[0, ~7]$ \\
$\rightarrow 3::1::[7, ~0, ~7]$ \\
$\rightarrow 3::[1, ~7, ~0, ~7]$ \\
$\rightarrow [3, ~1, ~7, ~0,~7]$ \\

\subsubsection{Informe de corrección}
Sea $L = \{ l~|~l~$ es un secuencia $<a_0,...a_i,...,a_n>$ para $0 \leq i \leq n \in  N^+ ~ \cup ~[\ ] \}$  un conjunto de secuencias de n elementos junto a un elemento que representa un lista vacia y sea $f: List[Int] \rightarrow List[Int]$ una funcińn definida en L para $k \geq 1$ número de elementos, y codominio en conjunto L, tal que $f(l) = m$ cumple $l \in L$ y $m \in L \land m \subseteq l \lor m = \emptyset$ y representa la secuencia de los elementos $a_i \in l$ que son menores a un $k$ determinado. \\
Sea $Pf$ el anterior programa realizado en $Scala$ que implementa $f$ y al cúal se quiere demostrar su correctitud:\\

Se quiere probar \\ \\
\begin{math}
\centering
\forall \in N \setminus \{0\} : P_f(List(a_1,...,a_i,...,a_n)) = f(List(a_1,...,a_i,...,a_n))
\end{math}
\\
\begin{itemize}
\item{Caso base: Cuando $l = [\ ]$. Utilizando el modelo de sustitución, se tiene:}

$P_f([\ ]) \rightarrow si(l = [\ ]) \ [\ ]$ \\ \\
Se tiene en la funcion teorica $f$
$f([]\ ) = [\ ]$\\ \\
$\therefore ~ P_f([\ ]) = f([\ ])$ \\
\item{Caso de inducción: para $n = k + 1$ se debe tener que }\\

\begin{math}
Si~P_f(List(b_1,...,b_i,...,b_k)) = f (List(b_1,...,b_i,...,b_k)) \\
\rightarrow P_f(List(b_1,...,b_i,...,b_{k + 1})) = f (List(b_1,...,b_i,...,b_{k + 1}))
\end{math}

Sea $m = List(a_1,...,a_i,..., a_{k +1})$ y p un valor comparar en el predicado $a_i \in m \land p \geq 1$, y por medio del modelo de sustitucion en $P_f(m)$ se tiene\\

$P_f(m) \rightarrow if~(m.isEmpty)~[\ ]~else~if( m.head < p)~m.head::\\~P_f(~List(a_2,...,a_i,..,a_{k+1}))~else~P_f(List(a_2,...,a_i,..,a_{k+1}))$ \\ \\
Sea $\alpha = List(a_2,...,a_i,...,a_{k+1})$ la lista resultante en el primer paso de la iteracion y diferenciando aritméticamente los subíndices de los elementos extremos de las lista, se tiene \\
$ k + 1 - 2 = k~-~1$ \\
Lo cual es tamaño obtenido para una lista de $k$ elementos, así, por hipotésis de indución, para una lista de k elementos conlleva a\\ \\
$\therefore P_f(List(a_1,..,a_i,...,a_n)) = f(List(a_1,..,a_i,...,a_n))$ $\forall n \in N \setminus \{0\}$

\end{itemize}



\subsection{Calculando el k-ésimo elemento de una lista}
\subsubsection{Informe de procesos}
\begin{lstlisting}[style=scalaStyle, caption=Calcula el k-ésimo elemento de una lista en el orden natural de los enteros positivos]
 def k_elem(l : List[Int], k : Int) : Int = {

    if(tamI(menoresQue(l, l.head)) == k - 1) l.head
    else if(tamI(menoresQue(l, l.head)) >= k) k_elem(menoresQue(l, l.head), k)
    else k_elem(noMenoresQue(l, l.head), k - tamI(menoresQue(l, l.head)) - 1)
  }
\end{lstlisting}
$k\_elem(List(4, ~6, ~1, ~2, ~7, ~10, ~8, ~5, ~3), ~6)$\\
$\rightarrow if(tamI(menoresQue(List(4, ~6, ~1, ~2, ~7, ~10, ~8, ~5, ~3)), ~4)) ~== ~6 ~- ~1) \\
~~~~~List(4, ~6, ~1, ~2, ~7, ~10, ~8, ~5, ~3).head\\
\ \ \ ~ ~~ else ~if(tamI(menoresQue(List(4, ~6, ~1, ~2, ~7, ~10, ~8, ~5, ~3), 4)) >= ~6)$
$\\ ~~~~~k\_elem(menoresQue(List(4, ~6, ~1, ~2, ~7, ~10, ~8, ~5, ~3), \\
~~~~~~List(4, ~6, ~1, ~2, ~7, ~10, ~8, ~5, ~3).head), ~6)$ \\
$\ ~~~else ~k\_elem(noMenoresQue( List(4, ~6, ~1, ~2, ~7, ~10, ~8, ~5, ~3), \\
~~~~~~~~List(4, ~6, ~1, ~2, ~7, ~10, ~8, ~5, ~3).head), \\
~~~~~~~~~~~~6 - tamI(menoresQue(List(4, ~6, ~1, ~2, ~7, ~10, ~8, ~5, ~3), \\
~~~~~~~~~~~~~List(4, ~6, ~1, ~2, ~7, ~10, ~8, ~5, ~3).head)) ~-~ 1)$ \\ \\
$\twoheadrightarrow ~k\_elem(noMenoresQue( List(4, ~6, ~7, ~10, ~8, ~5), \\
~~~~~~~~List(4, ~6, ~7, ~10, ~8, ~5).head),
~6 - tamI(menoresQue(List(4, ~6, ~7, ~10, ~8, ~5), \\
~~~~~~~~~~~~~List(4, ~6, ~7, ~10, ~8, ~5).head)) ~-~ 1)$ \\ \\
$\twoheadrightarrow ~k\_elem(noMenoresQue( List(4, ~6, ~7, ~10, ~8, ~5), \\
~~~~~~~~List(4, ~6, ~7, ~10, ~8, ~5).head),
~6 - 3 ~-~ 1)$
$~k\_elem(List(4, ~6, ~7, ~10, ~8, ~5), 2)$ \\ \\
$\rightarrow~k\_elem(List(~6, ~7, ~10, ~8, ~5), 1)$ \\
$~~~if(tamI(menoresQue(List(~6, ~7, ~10, ~8, ~5)), ~6)) ~== ~1 ~- ~1) \\
~~~~~List(~6, ~7, ~10, ~8, ~5).head\\
\ \ \ ~ ~~ else ~if(tamI(menoresQue(List( ~6, ~7, ~10, ~8, ~5), 4)) >= ~1)$
$\\ ~~~~~k\_elem(menoresQue(List(6, ~7, ~10, ~8, ~5),
~~List( ~6, ~7, ~10, ~8, ~5).head), ~1)$ \\
$\ ~~~else ~k\_elem(noMenoresQue( List(~6, ~7, ~10, ~8, ~5 ), \\
~~~~~~~~List(~6 ~7, ~10, ~8, ~5).head),
~1 - tamI(menoresQue(List( ~6, ~7, ~10, ~8, ~5 ), \\
~~~~~~~~List(~6, ~7, ~10, ~8, ~5).head)) ~-~ 1)$ \\
$\rightarrow ~~6$
\subsubsection{Informe de corrección}
\begin{itemize}
\item{$s_o = (L, k)$  L la lista original, v el valor pivote}
\item{$s = (L, k)$  }
\item{$s_f =~ k = f(L) \land tamI(menoresQue(L)) == k$ }
\item{$inv(s_i)$ = $ 0 \leq k \leq n \land k = f(L)$ }
\item{$T(s_i) = (L, k), m, v$  donde si $tamI(menoresQue(L)) \geq L = List(a_j,...,L.head)$ sino $L = List(L.head, a_m) para ciertos 0 \leq m, j \geq n \land k = k-1-tamI(menoresQue(L.tail, L.head))$}
  \end{itemize}

\begin{itemize}
\item{$Inv(s_0).$ Cumple con el invariante el estado inicial, $s_0 = (L, 0)$}
\item{$s_i \neq s_f \land Inv(s_i) \rightarrow Inv(T(s_i))$ El nuevo estado cumple con la condicion invariante }
$s_i \neq s_f \land Inv(s_i) \rightarrow Inv(T(s_i))$  \\
$\equiv$ \\
 $0\leq k \leq n \land k = f(L)$
 $\rightarrow k = k \lor K = k - tamI(menoresQue(l, l.head)) -1$
\item{$s_f =~ k = f(L) \land tamI(menoresQue(L)) == k$ }
\item{$inv(s_i)$ = $ 0 \leq k \leq n \land k = f(L)$ }
\item{$T(s_i) = (L, k), m, v$  donde si $tamI(menoresQue(L)) \geq L = List(a_j,...,L.head)$ sino $L = List(L.head, a_m) para ciertos 0 \leq m, j \geq n \land k = k-1-tamI(menoresQue(L.tail, L.head))$}
  \end{itemize}
\subsection{Ordenando una lista}
\subsubsection{Informe de procesos}
\subsubsection{Informe de corrección}
\section{Funciones de alto orden implementadas}

A continuación, se presenta la funciones implementadas de alto orden, las cuales fueron utilizadas para instanciar otras funciones (funciones generadoras), a través de su paso como parámetro, ya sea referenciada (nominada) o como funcion anónima(inline), o como valor retorno de la misma. \\[16pt]

\begin{table}
 \scriptsize
\begin{tabular}{ |p{4cm}||p{3cm}|p{5.5cm}|  }
 \hline
 \multicolumn{3}{|c|}{Funciones de alto orden} \\
 \hline
 Función& Forma de alto orden  & Expresión donde aparece\\
 \hline
  $Chip$ & Retorno & Retorno de funciones crearChipunario, crearChipBinario, half\_adder, full\_adder, adder\\
  \hline
  $(x:Int) \Rightarrow (x - 1)$  & Lambda como argumento &  $crearChipUnario((x:Int) \Rightarrow (x - 1)):Chip$\\
  \hline
  $ (x: Int, y: Int ) => (x * y)$  & Lambda como argumento &  $crearChipBinario((x: Int, y: Int ) => (x * y)): Chip$\\
  \hline
  $(x: Int, y: Int ) => (x + y) - (x * y )$  & Lambda como argumento &  $crearChipBinario((x: Int, y: Int ) => (x + y) - (x * y )): Chip$\\
\hline
  $half\_adder$ & Variable la cual se asigna una función de retorno & val $half\_adder$ = $( operands: List[Int]) => \{~...~ \}$\\
\hline
  $full\_adder$& Variable la cual se asigna una función de retorno  & val $full\_adder$ = $( operands: List[Int]) => \{~...~ \}$\\
\hline
  $adder$& Variable la cual se asigna una función de retorno & val $adder$ = $( operands: List[Int]) => \{~...~ \}$ \\
 \hline
\end{tabular}
\centering
\caption{Funciones de alto orden realizadas en la implementación del circuito lógico.}
\label{tabla:1}
\end{table}

\section{Crear chip unario}
\subsection{Informe de procesos}
Realiza una operación logica sobre un solo valor de entrada. A continuación, se presenta su implementación en $Scala$ \\[16pt]
\begin{lstlisting}[style=scalaStyle, caption=Aplica una operación binaria sobre una valor de entrada.]
  def crearChipUnario( f: Int => Int ) : Chip = (arg: List[Int]) => { // Apply the f function on the head of current list and call recursively the crearChipUnarioHelper with function f, a accumulated list with new transformed value as its head, and the current list tail, until the empty list condition is reached.
    @tailrec
    def crearChipUnarioHelper(f: Int => Int, transformedList: List[Int],  currentList: List[Int]): List[Int] =
      if (currentList.isEmpty) transformedList
      else crearChipUnarioHelper(f, f(currentList.head)::transformedList, currentList.tail)

    // The initial state of the iteration
    crearChipUnarioHelper(f, List(), arg)
  }
\end{lstlisting}
\subsection{Informe de corrección}
$val~chip\_not = crearchipUnario(x => 1- x)$ \\ \\
Caso 1: \\ 
$chip\_not((List(0)))$ \\
$\rightarrow crearChipUnarioHelper(x => 1 - x, [~], ~List(0))$ \\
$\rightarrow if (List(0).isEmpty)~[~] $ \\
$~~~~~else~crearChipUnarioHelper(x => 1 - x, (1 - 0)::[~], [~])$ \\
$\rightarrow if (List().isEmpty)~[1] $ \\
$~~~~~else~crearChipUnarioHelper(x => 1 - x, (1 - 1)::[1], [~])$ \\
$\rightarrow [1]$ \\ \\
Caso 2: \\ 
$chip\_not((List(1)))$ \\
$\rightarrow crearChipUnarioHelper(x => 1- x, [~], List(1))$ \\
$\rightarrow if (List(1).isEmpty)~[~] $ \\
$~~~~~else~crearChipUnarioHelper(x => 1 - x, (1 - 1)::[~], [~])$ \\
$\rightarrow if (List().isEmpty)~[0] $ \\
$~~~~~else~crearChipUnarioHelper(x => 1 - x, (1 - 0)::[0], [~])$ \\
$\rightarrow [0]$ \\ \\
\section{Conclusion}
La conclusion
\newpage{}
\begin{displaymath}
 a = \sum F \dot m = \frac{dv}{dt}
\end{displaymath}

\end{document}
List[Int] \Rightarrow List[Int]